% Importações de pacotes
\usepackage[T1]{fontenc}
\usepackage[utf8]{inputenc}
\usepackage[brazil]{babel}
\usepackage{xcolor}

\usepackage[utf8]{inputenc}                         % Acentuação direta
\usepackage[T1]{fontenc}                            % Codificação da fonte em 8 bits
\usepackage{graphicx}                               % Inserir figuras
\usepackage{amsfonts, amssymb, amsmath}             % Fonte e símbolos matemáticos
\usepackage{booktabs}                               % Comandos para tabelas
\usepackage{verbatim}                               % Texto é interpretado como escrito no documento
\usepackage{multirow, array}                        % Múltiplas linhas e colunas em tabelas
\usepackage{indentfirst}                            % Endenta o primeiro parágrafo de cada seção.
\usepackage{listings}                               % Utilizar codigo fonte no documento
\usepackage{xcolor}
\usepackage{microtype}                              % Para melhorias de justificação?
\usepackage[portuguese,ruled,lined]{algorithm2e}    % Escrever algoritmos
\usepackage{algorithmic}                            % Criar Algoritmos  
%\usepackage{float}                                 % Utilizado para criação de floats
\usepackage{amsgen}
\usepackage{lipsum}                                 % Usar a simulação de texto Lorem Ipsum
%\usepackage{titlesec}                              % Permite alterar os títulos do documento
\usepackage{tocloft}         % Permite alterar a formatação do Sumário
\usepackage{etoolbox}                               % Usado para alterar a fonte da Section no Sumário
\usepackage[nogroupskip,nonumberlist]{glossaries}   % Permite fazer o glossario

\usepackage[font={singlespacing},format=plain,skip=0pt,singlelinecheck = false]{caption}            % Altera o comportamento da tag caption

\usepackage[alf, abnt-emphasize=bf, recuo=0cm, abnt-etal-cite=3, abnt-etal-list=0, abnt-etal-text=it, bibjustif]{abntex2cite}  % Citações padrão ABNT
%\usepackage[bottom]{footmisc}                      % Mantém as notas de rodapé sempre na mesma posição
%\usepackage{times}                                 % Usa a fonte Times
%%%%%%%%%%%%%%%%%%% AVISO %%%%%%%%%%%%%%%%%%%%%%%%%%%%%%%%%%%%%%%%
%descomente as duas linhas abaixo para alterar o texto de Times New Roman para Arial:

%\usepackage{helvet}
%\renewcommand{\familydefault}{\sfdefault}  % Usa a fonte Arial              
%%%%%%%%%%%%%%%%%%%%%%%%%%%%%%%%%%%%%%%%%%%%%%%%%%%%%%%%%%%%%%%%%%

\usepackage{mathptmx}         % Usa a fonte Times New Roman			%\usepackage{lmodern}         % Usa a fonte Latin Modern
%\usepackage{subfig}          % Posicionamento de figuras
%\usepackage{scalefnt}        % Permite redimensionar tamanho da fonte
%\usepackage{color, colortbl} % Comandos de cores
%\usepackage{lscape}          % Permite páginas em modo "paisagem"
%\usepackage{ae, aecompl}     % Fontes de alta qualidade
%\usepackage{picinpar}        % Dispor imagens em parágrafos
%\usepackage{latexsym}        % Símbolos matemáticos
%\usepackage{upgreek}         % Fonte letras gregas
\usepackage{appendix}         % Gerar o apendice no final do documento
\usepackage{paracol}          % Criar paragrafos sem identacao
\usepackage{lib/ufersatex}	      % Biblioteca com as normas para trabalhos academicos
\usepackage{pdfpages}         % Incluir pdf no documento
\usepackage{amsmath}          % Usar equacoes matematicas

% importar meus pacotes
\usepackage{setspace}

\makeglossaries % Organiza e gera a lista de abreviaturas, simbolos e glossario
% use o comando \gls{} ao citar no texto para a sigla ficar visível na Lista. Caso a lista de abreviaturas não estiver aparecendo é porque você não utilizou o comando "gls" no texto. Exemplo no texto:   De acordo com o \gls{MCEG} ...  
\newacronym{PDF}{PDF}{\textit{Portable Document Format}}
\newacronym{API}{API}{\textit{Application Programming Interface}}
\newacronym{DQN}{DQN}{\textit{Deep Q-Network}}
\newacronym{EV}{EV}{\textit{Expected Value}}
\newacronym{GPU}{GPU}{\textit{Graphics Processing Unit}}
\newacronym{HTML}{HTML}{\textit{HyperText Markup Language}}
\newacronym{HTTP}{HTTP}{\textit{Hypertext Transfer Protocol}}
\newacronym{IA}{IA}{Inteligência Artificial}
\newacronym{JSON}{JSON}{\textit{JavaScript Object Notation}}
\newacronym{MCCFR}{MCCFR}{\textit{Monte Carlo Counterfactual Regret Minimization}}
\newacronym{MCTS}{MCTS}{\textit{Monte Carlo Tree Search}}
\newacronym{MVC}{MVC}{\textit{Model-View-Controller}}
\newacronym{PARC}{PARC}{\textit{Palo Alto Research Center}}
\newacronym{ReLU}{ReLU}{\textit{Rectified Linear Unit}}
\newacronym{RNA}{RNA}{Rede Neural Artificial}
\newacronym{SIGAA}{SIGAA}{Sistema Integrado de Gestão de Atividades Acadêmicas}
\newacronym{SISBI}{SISBI}{Sistema de Bibliotecas da UFERSA}
\newacronym{SUTIC}{SUTIC}{Superintendência de Tecnologia da Informação e Comunicação}
\newacronym{TCC}{TCC}{Trabalho de Conclusão de Curso}
\newacronym{UFERSA}{UFERSA}{Universidade Federal Rural do Semi-Árido}
\newacronym{USP}{USP}{Universidade de São Paulo} % Carrega as definições de siglas e abreviaturas
\makeindex      % Gera o Indice do documento

\usepackage{chngcntr}
\counterwithout*{footnote}{chapter}

\usepackage{listings}
\usepackage{xcolor} % Opcional, para melhorar as cores

% Configuração básica para o código não vazar da página e ter números nas linhas
\lstset{
    basicstyle=\ttfamily\footnotesize,
    breaklines=true,
    frame=single,
    numbers=left,
    numberstyle=\tiny\color{gray}
}

\usepackage{color}

% Cartas de pôquer em estilo simples e estável para uso inline
\newcommand{\pokercard}[3]{%
    {\setlength{\fboxsep}{0.35ex}%
        \textcolor{#2}{\fbox{\sffamily\bfseries #1\ensuremath{#3}}}%
    }%
}
\newcommand{\As}{\pokercard{A}{black}{\spadesuit}}
\newcommand{\Ah}{\pokercard{A}{red}{\heartsuit}}
\newcommand{\Ad}{\pokercard{A}{red}{\diamondsuit}}
\newcommand{\Ac}{\pokercard{A}{black}{\clubsuit}}
\newcommand{\Ks}{\pokercard{K}{black}{\spadesuit}}
\newcommand{\Kh}{\pokercard{K}{red}{\heartsuit}}
\newcommand{\Kd}{\pokercard{K}{red}{\diamondsuit}}
\newcommand{\Kc}{\pokercard{K}{black}{\clubsuit}}
\newcommand{\Qs}{\pokercard{Q}{black}{\spadesuit}}
\newcommand{\Qh}{\pokercard{Q}{red}{\heartsuit}}
\newcommand{\Qd}{\pokercard{Q}{red}{\diamondsuit}}
\newcommand{\Qc}{\pokercard{Q}{black}{\clubsuit}}
\newcommand{\Js}{\pokercard{J}{black}{\spadesuit}}
\newcommand{\Jh}{\pokercard{J}{red}{\heartsuit}}
\newcommand{\Jd}{\pokercard{J}{red}{\diamondsuit}}
\newcommand{\Jc}{\pokercard{J}{black}{\clubsuit}}
\newcommand{\tens}{\pokercard{10}{black}{\spadesuit}}
\newcommand{\tenh}{\pokercard{10}{red}{\heartsuit}}
\newcommand{\tend}{\pokercard{10}{red}{\diamondsuit}}
\newcommand{\tenc}{\pokercard{10}{black}{\clubsuit}}
\newcommand{\nines}{\pokercard{9}{black}{\spadesuit}}
\newcommand{\nineh}{\pokercard{9}{red}{\heartsuit}}
\newcommand{\nined}{\pokercard{9}{red}{\diamondsuit}}
\newcommand{\ninec}{\pokercard{9}{black}{\clubsuit}}
\newcommand{\eigs}{\pokercard{8}{black}{\spadesuit}}
\newcommand{\eigh}{\pokercard{8}{red}{\heartsuit}}
\newcommand{\eigd}{\pokercard{8}{red}{\diamondsuit}}
\newcommand{\eigc}{\pokercard{8}{black}{\clubsuit}}
\newcommand{\sevs}{\pokercard{7}{black}{\spadesuit}}
\newcommand{\sevh}{\pokercard{7}{red}{\heartsuit}}
\newcommand{\sevd}{\pokercard{7}{red}{\diamondsuit}}
\newcommand{\sevc}{\pokercard{7}{black}{\clubsuit}}
\newcommand{\sixs}{\pokercard{6}{black}{\spadesuit}}
\newcommand{\sixh}{\pokercard{6}{red}{\heartsuit}}
\newcommand{\sixd}{\pokercard{6}{red}{\diamondsuit}}
\newcommand{\sixc}{\pokercard{6}{black}{\clubsuit}}
\newcommand{\fives}{\pokercard{5}{black}{\spadesuit}}
\newcommand{\fiveh}{\pokercard{5}{red}{\heartsuit}}
\newcommand{\fived}{\pokercard{5}{red}{\diamondsuit}}
\newcommand{\fivec}{\pokercard{5}{black}{\clubsuit}}
\newcommand{\fours}{\pokercard{4}{black}{\spadesuit}}
\newcommand{\fourh}{\pokercard{4}{red}{\heartsuit}}
\newcommand{\fourd}{\pokercard{4}{red}{\diamondsuit}}
\newcommand{\fourc}{\pokercard{4}{black}{\clubsuit}}
\newcommand{\tres}{\pokercard{3}{black}{\spadesuit}}
\newcommand{\treh}{\pokercard{3}{red}{\heartsuit}}
\newcommand{\tred}{\pokercard{3}{red}{\diamondsuit}}
\newcommand{\trec}{\pokercard{3}{black}{\clubsuit}}
\newcommand{\twos}{\pokercard{2}{black}{\spadesuit}}
\newcommand{\twoh}{\pokercard{2}{red}{\heartsuit}}
\newcommand{\twod}{\pokercard{2}{red}{\diamondsuit}}
\newcommand{\twoc}{\pokercard{2}{black}{\clubsuit}}

% Definindo algumas cores bonitas para o código
\definecolor{lightgray}{rgb}{0.95, 0.95, 0.95}
\definecolor{darkgray}{rgb}{0.4, 0.4, 0.4}
\definecolor{purple}{rgb}{0.65, 0.12, 0.82}
\definecolor{blue}{rgb}{0.13, 0.13, 1}
\definecolor{green}{rgb}{0, 0.5, 0}

% Criando a definição do JavaScript
\lstdefinelanguage{JavaScript}{
    keywords={typeof, new, true, false, catch, function, return, null, catch, switch, var, let, const, if, in, while, do, else, case, break},
    keywordstyle=\color{blue}\bfseries,
    ndkeywords={class, export, boolean, throw, implements, import, this},
    ndkeywordstyle=\color{darkgray}\bfseries,
    identifierstyle=\color{black},
    sensitive=false,
    comment=[l]{//},
    morecomment=[s]{/*}{*/},
    commentstyle=\color{green}\ttfamily,
    stringstyle=\color{purple}\ttfamily,
    morestring=[b]',
    morestring=[b]"
}

% Configurando o padrão do listings para usar o que acabamos de criar
\lstset{
    language=JavaScript,
    backgroundcolor=\color{lightgray},
    extendedchars=true,
    basicstyle=\footnotesize\ttfamily,
    showstringspaces=false,
    showspaces=false,
    numbers=left,
    numberstyle=\tiny\color{gray},
    numbersep=9pt,
    tabsize=2,
    breaklines=true,
    showtabs=false,
    captionpos=b
}

\lstset{
    basicstyle=\ttfamily\footnotesize,
    breaklines=true,
    frame=single,
    showstringspaces=false,
    numbers=none,
    backgroundcolor=\color{gray!10},
    % --- SOLUÇÃO DOS ACENTOS AQUI ---
    literate=
        {á}{{\'a}}1 {é}{{\'e}}1 {í}{{\'i}}1 {ó}{{\'o}}1 {ú}{{\'u}}1
        {Á}{{\'A}}1 {É}{{\'E}}1 {Í}{{\'I}}1 {Ó}{{\'O}}1 {Ú}{{\'U}}1
        {à}{{\`a}}1 {è}{{\`e}}1 {ì}{{\`i}}1 {ò}{{\`o}}1 {ù}{{\`u}}1
        {À}{{\`A}}1 {È}{{\'E}}1 {Ì}{{\`I}}1 {Ò}{{\`O}}1 {Ù}{{\`U}}1
        {â}{{\^a}}1 {ê}{{\^e}}1 {î}{{\^i}}1 {ô}{{\^o}}1 {û}{{\^u}}1
        {Â}{{\^A}}1 {Ê}{{\^E}}1 {Î}{{\^I}}1 {Ô}{{\^O}}1 {Û}{{\^U}}1
        {ã}{{\~a}}1 {õ}{{\~o}}1 {Ã}{{\~A}}1 {Õ}{{\~O}}1
        {ç}{{\c c}}1 {Ç}{{\c C}}1
}